% ==========================================
% F2 - THU LOI RIP TREN L1
% Nhom K = 7126
% ==========================================

\section{F2 -- Thử nghiệm hội tụ RIP}
\label{sec:F2}

\subsection{Mục tiêu}

Thử nghiệm F2 kiểm tra thời gian A1 bỏ các tuyến RIP qua A2
khi liên kết L1 chỉ bị tắt ở đầu A2, đồng thời theo dõi ping
liên tục giữa LAN-B và LAN-A. Kết quả được đối chiếu với
các timer RIP để giải thích vì sao A1 không loại bỏ tuyến ngay
khi cổng phía A2 bị shutdown.

\subsection{Phương pháp thực hiện}

\subsubsection{Chuẩn bị và theo dõi kết nối}

Mô hình sử dụng RIPv2 trong AS A, gồm A1, A2 và A3.
Liên kết L1 nối FastEthernet0/0 của A1
(\texttt{10.71.26.1/30}) với FastEthernet0/0 của A2
(\texttt{10.71.26.2/30}). Đường thay thế giữa hai router
có thể đi qua A3 trên L2 và L3.

Hai LAN dùng để kiểm thử là LAN-B
\texttt{10.171.126.64/26}, đại diện bởi Loopback1 của B2
(\texttt{10.171.126.65}), và LAN-A
\texttt{10.171.126.0/26}, đại diện bởi Loopback1 của A1
(\texttt{10.171.126.1}). Trên B2, nhóm chạy lệnh:

\begin{lstlisting}
B2-7126#ping 10.171.126.1 source Loopback1 repeat 2000 timeout 1
\end{lstlisting}

Lệnh gửi tối đa 2000 gói ICMP Echo Request, kích thước
100 byte, với thời gian chờ mỗi phản hồi là 1 giây.
Địa chỉ nguồn thuộc LAN-B, nên phép thử kiểm tra kết nối
giữa hai LAN đúng yêu cầu Phần F.

Trên A1, nhóm bật timestamp và debug để theo dõi cập nhật RIP
cùng các thay đổi trong bảng định tuyến:

\begin{lstlisting}
A1-7126#configure terminal
A1-7126(config)#service timestamps debug uptime
A1-7126(config)#service timestamps log uptime
A1-7126(config)#logging console debugging
A1-7126(config)#end
A1-7126#debug ip routing
A1-7126#debug ip rip
\end{lstlisting}

\subsubsection{Gây lỗi trên L1}

Trong khi ping đang chạy, nhóm chỉ tắt cổng phía A2:

\begin{lstlisting}
A2-7126#configure terminal
A2-7126(config)#interface FastEthernet0/0
A2-7126(config-if)#shutdown
\end{lstlisting}

Ảnh minh chứng ghi nhận cổng trên A2 chuyển sang
trạng thái administratively down lúc \texttt{00:31:17.083}.
Line protocol chuyển down lúc \texttt{00:31:18.083}.
Cổng phía A1 không bị shutdown.
Theo đặc điểm Dynamips nêu trong đề bài,
đầu bên kia không tự thấy cổng bị tắt; vì vậy A1 phải
phát hiện việc mất cập nhật RIP qua L1.

\partzeroimage
{images/testing/F2/one_end_failure_start.png}
{Thử lỗi F2: shutdown L1 phía A2, bật debug trên A1
và ping từ LAN-B tới LAN-A; đồng hồ hiển thị 00:03:25.845}
{fig:F2-start}

\subsection{Kết quả kiểm thử và giải thích}

\subsubsection{A1 loại bỏ các tuyến qua A2}

Trong Hình~\ref{fig:F2-start}, A1 còn nhận cập nhật qua L1
từ A2 (\texttt{10.71.26.2}) lúc \texttt{00:31:10}.
Sau khi A2 tắt L1, A1 vẫn nhận cập nhật qua L2
từ A3 (\texttt{10.71.26.6}); ảnh thứ hai ghi nhận
hoạt động này lúc \texttt{00:33:53}.

\noindent\begin{minipage}{\linewidth}
Đến \texttt{00:34:16} trên A1, log xuất hiện các dòng
sau, trích từ ảnh minh chứng:

\begin{lstlisting}
00:34:16: RT: del 10.71.26.8/30 via 10.71.26.2, rip metric [120/1]
00:34:16: RT: NET-RED 10.71.26.8/30
00:34:16: RT: delete route to 10.71.26.12 via 10.71.26.2, rip metric [120/1]
00:34:16: RT: no routes to 10.71.26.12, entering holddown
00:34:16: RT: NET-RED 10.71.26.12/30
00:34:16: RT: delete route to 10.171.26.2 via 10.71.26.2, rip metric [120/1]
00:34:16: RT: no routes to 10.171.26.2, entering holddown
00:34:16: RT: NET-RED 10.171.26.2/32
\end{lstlisting}
\end{minipage}

Các dòng \texttt{del} và \texttt{delete route} xác nhận
A1 bỏ đường qua next-hop A2 \texttt{10.71.26.2}.
Đối với mạng L4 \texttt{10.71.26.12/30}
và Loopback0 của A2 \texttt{10.171.26.2/32},
log còn ghi nhận trạng thái \texttt{entering holddown}.


\partzeroimage
{images/testing/F2/one_end_routes_removed.png}
{A1 bỏ các tuyến RIP qua A2 tại uptime 00:34:16;
đồng hồ của bản ghi hiển thị 00:06:25.554}
{fig:F2-routes-removed}

\subsubsection{Ước lượng thời gian A1 bỏ tuyến}

Nhóm sử dụng đồng hồ hiển thị ở góc trên bên trái
của hai ảnh để ước lượng khoảng chờ:

\begin{table}[H]
\centering
\caption{Các mốc quan sát và kết quả đo trong thử nghiệm F2}
\label{tab:F2-results}
\small
\renewcommand{\arraystretch}{1.3}
\begin{tabularx}{\textwidth}{X c}
\toprule
\textbf{Nội dung} & \textbf{Giá trị} \\
\midrule
Khung hình đã ghi nhận shutdown trên A2 & 00:03:25.845 \\
Khung hình có log A1 bỏ tuyến qua A2 & 00:06:25.554 \\
Khoảng thời gian giữa hai khung hình & 179,709 giây \\
Timer Invalid để đối chiếu & 180 giây \\
\bottomrule
\end{tabularx}
\end{table}

\begin{equation}
\begin{aligned}
T_{\text{bỏ tuyến}}
&\approx (6 \times 60 + 25.554) - (3 \times 60 + 25.845)\\
&= \boxed{179.709\ \text{s}} \approx 180\ \text{s}.
\end{aligned}
\end{equation}

Giá trị trên là ước lượng từ hai khung hình được cung cấp,
không phải phép đo chính xác tới mili giây của thời điểm
thực thi lệnh và xử lý tuyến. Các timestamp uptime
trên A1 và timestamp hệ thống trên A2 thuộc hai router
khác nhau, nên không lấy chúng trừ trực tiếp cho nhau.

\subsubsection{Đối chiếu với các timer RIP}

Output \texttt{show ip protocols} trên A2 tại
Hình~\ref{fig:partB-a2-protocols} ghi nhận bộ timer
\texttt{30/180/180/240} giây. Ý nghĩa các timer được
đối chiếu với tài liệu lệnh \texttt{timers basic} của Cisco.\footnote{Cisco,
\href{https://www.cisco.com/c/en/us/td/docs/ios/iproute_rip/command/reference/irr_book/irr_rip.html}
{Cisco IOS IP Routing: RIP Command Reference -- timers basic (RIP)}.}

\begin{table}[H]
\centering
\caption{Các timer RIP dùng để đối chiếu thử nghiệm F2}
\label{tab:F2-timers}
\small
\renewcommand{\arraystretch}{1.3}
\begin{tabularx}{\textwidth}{l c X}
\toprule
\textbf{Timer} & \textbf{Giá trị (s)} & \textbf{Ý nghĩa} \\
\midrule
Update & 30 & Chu kỳ gửi cập nhật định kỳ. \\
Invalid & 180 & Ngưỡng không được làm mới để tuyến bị coi là không hợp lệ. \\
Holddown & 180 & Hạn chế tiếp nhận thông tin tuyến trong giai đoạn giữ tuyến. \\
Flush & 240 & Mốc loại bỏ thông tin tuyến cũ, tính từ lần cập nhật cuối. \\
\bottomrule
\end{tabularx}
\end{table}

Khoảng chờ gần 180 giây phù hợp với cơ chế hết hạn
Invalid khi A1 không còn nhận cập nhật qua L1.
Timer này tính từ lần làm mới tuyến cuối, nên thời gian
từ lúc shutdown không nhất thiết đúng 180 giây.
Log \texttt{entering holddown} phản ánh giai đoạn
tuyến mất hiệu lực; không dùng kết quả này để kết luận
Flush 240 giây đã hết hạn.

\subsubsection{Kết quả theo dõi ping giữa hai LAN}

Cửa sổ B2 trong hai ảnh hiển thị các dấu \texttt{!},
cho thấy đã có các phản hồi ICMP thành công.
Ảnh đầu còn cho thấy A1 đã có default route
\texttt{0.0.0.0/0} qua A3 \texttt{10.71.26.6}
trên FastEthernet0/1.
Vì vậy, việc A1 giữ một số tuyến nội bộ qua A2
không đồng nghĩa toàn bộ kết nối giữa hai LAN bị mất
trong suốt khoảng chờ bỏ tuyến. Trong suốt quá trình ping, không có hiện tượng gián đoạn gây mất gói


\subsection{Nhận xét và khôi phục cấu hình}

Thử nghiệm cho thấy khi L1 chỉ bị tắt ở đầu A2,
A1 không bỏ ngay các tuyến RIP học qua liên kết này.
Sau khoảng 180 giây theo hai khung hình,
debug ghi nhận A1 loại bỏ đường qua A2 và đưa
một số tuyến vào holddown, phù hợp với timer Invalid.

Để kết thúc thử nghiệm và chuẩn bị cho F3, cần bật lại
cổng trên A2, tắt debug trên A1 và chờ mạng ổn định:

\begin{lstlisting}
! A2-7126
configure terminal
interface FastEthernet0/0
 no shutdown
end

! A1-7126
undebug all
show ip route rip
show ip protocols
ping 10.171.26.2
\end{lstlisting}


\newpage
